\subsection{The smallest algebra and iteration}
\label{subsec:clifford-smallest-algebra}

Write $q(y)=|y|^2$ on the target.

\begin{theorem}[The smallest algebra under Clifford pullback]
\label{thm:smallest-clifford-pullback}
Let $P:\R^N\to\R^m$ be the Clifford map above and let $f$ be a
nonzero essential homogeneous polynomial of degree $s\ge2$ on
$\R^m$, so $V_{\mathrm{cyl}}(f)=\{0\}$. Then
\begin{equation}\label{eq:smallest-clifford-pullback}
 \cA_R(f\circ P)=\R[R,P^*\cA_q(f)].
\end{equation}
If $Q<R^2$ somewhere, this is isomorphic to the polynomial extension
$\cA_q(f)[t]$, with target degrees doubled and $\deg t=2$.
In particular, its minimal number of homogeneous generators is one
more than that of $\cA_q(f)$.
\end{theorem}

\begin{proof}
Set $\cA=\cA_q(f)$ and $\cB=\cA_R(f\circ P)$.
By \cref{lem:Clifford-algebra-pullback}, $\R[R,P^*\cA]$ is Laplacian
and contains $f\circ P$, so $\cB\subseteq\R[R,P^*\cA]$.
For the reverse inclusion put \(\cC=\{u\in\R[y_1,\ldots,y_m]:u\circ P\in\cB\}\).
This is a unital algebra containing $f$. If $u\in\cC$, the chain
rule gives $4R(\Delta u)\circ P=\Delta(u\circ P)\in\cB$.
The quotient is a polynomial, so
\cref{prop:MR-finite-and-division} gives $(\Delta u)\circ P\in\cB$.
Thus $\cC$ is closed under $\Delta$ and, by the product rule, $\Gamma$.

It remains to recover $q$. The recursion
\eqref{eq:MR-higher-contraction} uses only multiplication and
$\Delta$, so $B_{s-1}(f,f)\in\cC$. This quadratic form is
$\langle My,y\rangle$ for a positive semidefinite $M$ whose kernel
is exactly the cylindrical directions of $f$: by polarization,
$B_{s-1}(f,f)(v)=0$ if and only if $\partial_vf\equiv0$.
Essentiality of $f$ therefore makes $M$ positive definite. Since
\[
 \Gamma(\langle My,y\rangle,\langle M^jy,y\rangle)
       =4\langle M^{j+1}y,y\rangle,
\]
all positive powers of $M$ give quadratic forms in $\cC$.
Choose a polynomial $p$ with $p(0)=0$ and $p(\lambda)=1$ at every
eigenvalue of $M$. Then $p(M)=\Id$ gives $q\in\cC$.
Consequently $\cA\subseteq\cC$, proving
\eqref{eq:smallest-clifford-pullback}.

At a point where $Q<R^2$, the Gram matrix of
$\nabla R,\nabla P_1,\ldots,\nabla P_m$ has determinant \(4^{m+1}R^{m-1}(R^2-Q)>0\).
Thus $(R,P)$ has an open set in its
image, and $R$ is algebraically independent of the pulled-back
polynomial ring. This proves the asserted isomorphism. For a finitely
generated connected graded algebra, the minimal homogeneous generator
number is $\dim(\cA_+/\cA_+^2)$, where $\cA_+$ is its positive-degree
ideal. Adjoining $t$ adds exactly one independent class.
\end{proof}

When $Q=R^2$ identically, the relation $q(P)=R^2$ must be retained;
the conclusion about one additional generator does not apply.

Choose minimal homogeneous generators
$\cA_q(f)=\R[u_1,\ldots,u_r]$, with $\deg u_i=d_i$, and write
\[
 \Delta u_i=b_i(u),\qquad \Gamma(u_i,u_j)=c_{ij}(u),\qquad q=q_*(u).
\]
The lifted variables $U_i=u_i(P)$ have the differential table
\begin{equation}\label{eq:clifford-minimal-table}
\begin{gathered}
 \Delta R=2N,\qquad \Gamma(R,R)=4R,\qquad
 \Delta U_i=4R\,b_i(U),\\
 \Gamma(R,U_i)=4d_iU_i,\qquad
 \Gamma(U_i,U_j)=4R\,c_{ij}(U).
\end{gathered}
\end{equation}
Thus the target table determines the calculation for any potential
$G=R^a\phi(U)$; only $\Delta R$ depends on the source dimension.
For a target potential factoring through $\cA_q(f)$ as $H=\phi(u)$,
this applies when $\phi$ is $C^3$ on an open neighborhood of
$u(\R^m\setminus Z)$.

\begin{proposition}[Target equivalence and dimension monotonicity]
\label{prop:clifford-transfer-monotonicity}
Fix a Clifford map $P_0:\R^{N_0}\to\R^m$ and, for $k\ge2$, set
\[
 P_k(x_1,\ldots,x_k)=\sum_{\alpha=1}^kP_0(x_\alpha),
 \qquad N_k=kN_0,\qquad R=\sum_{\alpha=1}^k|x_\alpha|^2.
\]
Assume the target hypotheses on $f,Z,H,M$ of
\cref{thm:clifford-area-transfer}, except for $H\cL(H)\ge0$.
For fixed $a>0$, put $G_k=R^aH(P_k)$. Then
\[
 H\cL(H)\ge0\text{ off }Z
 \quad\Longleftrightarrow\quad
 G_k\cL(G_k)\ge0\text{ off }P_k^{-1}(Z)
 \text{ for some }k\ge2.
\]
Once the source inequality holds for one such $k$, it holds for every
larger $k$. In particular, the target inequality is equivalent to the
source inequalities in all sufficiently large members of this family.
\end{proposition}

\begin{proof}
The image of $(R,P_k)$ is exactly \(\{(z,y)\in\R\times\R^m:z\ge|y|\}\).
Indeed, $|P_k|\le R$. Conversely, choose a unit vector $e$ with
$y=|y|e$ (arbitrary if $y=0$), and write
$y=y_1+y_2$, where
$y_1=(z+|y|)e/2$ and $y_2=-(z-|y|)e/2$.
The eigenvector construction in
\cref{lem:Clifford-algebra-pullback} realizes $P_0(x_i)=y_i$
with $|x_i|^2=|y_i|$. Two blocks therefore suffice.

For $y\notin Z$ and $z\ge|y|$, put $b=a(a+2\ell)$ and
\[
 \Psi(z,y)=z^aH(y),\qquad
 \mathcal W(z,y)=4z^{2a+1}W_H(y)+4bz^{2a-1}H(y)^2.
\]
For $N=N_k$, let $\mathcal D_N(z,y)$ be the value of
$G_k\cL(G_k)$ at any point with $(R,P_k)=(z,y)$.
Formula \eqref{eq:clifford-transfer-defect} shows that this value
is well-defined, and \eqref{eq:clifford-transfer-gradient} gives
$|\nabla G_k|^2=\mathcal W(z,y)>0$ there.
For two block-sum dimensions $N'=N_{k'}\ge N=N_k$, the same
formula gives the finite difference
\begin{equation}\label{eq:clifford-dimension-monotonicity}
 \mathcal D_{N'}(z,y)-\mathcal D_N(z,y)
 =\frac{2a(N'-N)}{z}\Psi(z,y)^2\mathcal W(z,y)\ge0.
\end{equation}
The common image therefore proves monotonicity of the source
inequalities.

The target inequality implies the source inequality for large $k$
by \cref{thm:clifford-area-transfer}. Conversely, fix $y\notin Z$.
The image calculation permits $(R,P_k)=(z,y)$ for every $z\ge|y|$.
By \eqref{eq:clifford-transfer-defect},
\[
 \lim_{z\to\infty}
 \frac{\mathcal D_{N_k}(z,y)}{16z^{4a+2}}
     =H(y)\cL(H)(y).
\]
If the source inequality holds, each expression on the left is
nonnegative. This proves necessity.
\end{proof}

Composition preserves umbilic maps
\cite[Proposition~5.1.8(1)]{RiedlerThesis25}. Induction from
\cref{thm:clifford-area-transfer,thm:smallest-clifford-pullback}
gives the following conclusions.

\begin{corollary}[Iteration]
\label{cor:iterated-clifford-transfer}
Let $f$ be a nonzero essential homogeneous polynomial of degree
$s\ge2$ on the target $\R^m$, and suppose $f,Z,H,M$ satisfy
\cref{thm:clifford-area-transfer}, with $H=\phi(u)$ as above.
Successive Clifford block sums with at least two blocks and
sufficiently large source dimensions yield minimizing phase boundaries.
After $j$ steps the defining degree is $2^js$ and the smallest algebra
is isomorphic to $\cA_q(f)[t_1,\ldots,t_j]$, with minimal generator
degrees
\[
 2^jd_1,\ldots,2^jd_r,\qquad 2,4,\ldots,2^j.
\]
Thus its minimal homogeneous generator number is exactly $r+j$.
The same ratio constant $M$ can be used at every step. If both signs
of $f$ occur near every zero, the minimizing supports are the complete
successive zero sets.
\end{corollary}

\begin{proof}
The ratio estimate \eqref{eq:clifford-transfer-ratio} preserves $M$.
The hypotheses on regularity, signs, the exceptional set, and the
nonvanishing of the gradient persist by the proof of \cref{thm:clifford-area-transfer};
the new potential remains a function of the new generators.
If $\partial_v(g\circ P)\equiv0$ for a nonzero homogeneous $g$
of degree $d>0$, then $v=0$: otherwise
$(g\circ P)(v+x)=(g\circ P)(x)$ gives vanishing order $2d$ at $v$,
whereas the submersion property of $P$ makes this equal to the
vanishing order of $g$ at $P(v)$, which is at most $d$.
The common-image calculation above gives $Q<R^2$ somewhere at each
step, so the algebra and degree assertions follow inductively from
\cref{thm:smallest-clifford-pullback}. At each stage a symmetric
Clifford system exists on $\Lambda^*\R^m$, with generators given by
exterior multiplication plus contraction; direct sums make its source
dimension arbitrarily large. Thus each required dimension bound can
be met. The support assertion follows at every stage from
\cref{thm:clifford-area-transfer}.
\end{proof}
